
\section*{Alban pendant les guerres de religion (suite)}

\lettrine{S}{i} octobre 1571 a vu le départ des Protestants de notre cité, les hostilités n'ont pas fini pour autant. Des bandes de pillards, profitant des querelles religieuses, sillonnent le pays. On pille et on massacre. Quand on attaque un village catholique, on se prétend Huguenot, et pour piller une localité protestante, on se dit catholique. Nos braves gens ne se méprennent point sur la valeur de ces convictions opportunistes et finissent par désigner ces combattants d'un même titre : « les voleurs ». Sous la conduite de Bertrand de Rabastens, vicomte de Paulin, ils s'installent à Alban où ils se trouvent encore en 1579. À cette époque, le duc de Montmorency, gouverneur de Languedoc, ordonne de raser : « Alban, Villeneuve, Trébas… » qui, dit-il, sont « des repaires de voleurs ». En 1580, au mois de janvier, Alban, de nouveau aux mains des catholiques, retrouve sa tranquillité. Son seigneur, cependant, Jean Albert, a été emmené en captivité. La population est en telle détresse qu'elle ne peut payer les impôts, « les tailles » comme on dit alors, et comme encore on le dit en notre occitan ou patois local. En mai 1581, les « voleurs » sont de nouveau de retour, mais c'est pour peu de temps, car en juin de la même année, le gouverneur de Languedoc, écrivant d'Agde à l'évêque d'Albi, lui dit : « Veillez à la conservation d'Alban qu'on a voulu surprendre récemment ». Les catholiques sont donc maîtres de la place, mais tellement ruinés qu'ils maltraitent le porteur des rôles d'imposition. Il en est de même des huguenots qui, n'ayant pu tenir chez nous, se sont rabattus sur le Fraussé et s'y sont installés. Ils se refusent eux aussi à payer des impôts « se jactant » (se vantant) de traiter ceux qui viendraient les réclamer de telle sorte qu'ils n'auraient pas envie d'y retourner. Tandis que Jean Albert, seigneur d'Alban, s'alanguissait dans les prisons de Castres où les protestants l'ont enfermé, c'est un autre Albanais, un marchand lui aussi, qui, en 1581 devient le fermier de la « Claverie d'Alban ».
